\chapter[Sermon 67. The Angels of Seven Plagues Are Brought Forth. Moreover the Triumph and Praise of Christ's Holy Martyrs Is Described.]{The Angels of Seven Plagues Are Brought Forth. Moreover the Triumph and Praise of Christ's Holy Martyrs Is Described.}
\label{sermon:067}
\markright{Sermon 67. The Angels of Seven Plagues Are Brought Forth ...}

And I saw another sign in heaven, great and wonderful. Seven angels, having the seven final plagues. For in them is fulfilled the wrath of God. And I saw as it were a glassy sea mingled with fire, and those who had conquered the beast, and his image, and his mark, and the number of his name, standing on the glassy sea, having the harps of God: and they sang the song of Moses, the servant of God, and the song of the Lamb, saying: Great and marvelous are your works, Lord God Almighty, just and true are your ways, O King of Saints. Who shall not fear you and glorify your name? For you alone are holy; for all nations shall come and worship before you, for your judgments are made manifest.

\index{Antichrist}\index{Rome}Concerning the harvest and vintage, explained in the last part of the previous vision, is appended the fifth part of this godly work, which presents to us the fourth vision of this work, which some consider to be the fifth. This concerns God’s judgments and has two parts; therefore, it might also be divided into more visions, but we prefer to use fewer. First, he discusses at length the pains or torments prepared by God and to be executed upon Antichrist’s followers, and all the ungodly. Here is treated the judgment of the whore of Babylon, the destinies and ruin of Rome and the Roman church, the rejoicing and song of the saints, the coming of the judge for judgment, and the pain and everlasting destruction of all the wicked. These matters are addressed in chapters 15, 16, 17, 18, 19, and 20. Then also he reasons excellently of the reward of the saints and of everlasting blessedness throughout chapter 21. Everywhere, hell itself and heaven itself are set open; and we are given a glimpse, while in this mortal flesh, even into hell itself and into the very palace of Heaven. You will not find anywhere in all the Scriptures a continuous and abundant discussion of the judgments of God, the torments of the wicked, and the felicity and joys of the godly, as in this present work.

\index{John, Saint}And this treatise is especially necessary in this our last and ungracious world, wherein men, neglecting the spirit of God, have become like brute beasts, altogether carnal, regarding the flesh and wholly depending on it. Happy are all the victorious, wealthy, honorable, and glorious Antichristians; miserable are the poor and despised true Christians, and subject to the injuries and persecutions of all men. Therefore, carnal men esteem all things of the present fortune, and proclaim that their religion and conversation pleases God, and that the Christian way displeases. The godly are here also grievously tempted, as they were also in times past, read Psalm 73 and the first chapter of Habakkuk. The ungodly promise themselves that they shall reign forever; at length also they contain the judgments of God, neither thinking that it will ever come to pass that they shall be punished. The talk of punishments is devised by melancholic persons, and uttered of malice; and therefore they say and think them not to be regarded, but to be merry in this world. Therefore it behooved the place of God’s judgments to be most largely and diligently decided, and to be set as it were before the eyes of the hearers: to the end all might rightly understand what should be assuredly the end of good and evil. But the punishments of the ungodly are diverse, to wit, of this life present, and to come. And the punishments of this present life are almost innumerable; and the torments of the life to come are eternal and unspeakable; and as there is no comparison between the pained and true fire, so is there none between the punishments of this present life and that to come. But if men would earnestly believe that unspeakable joys and everlasting torments are prepared by God for good and evil, doubtless all would sin less and serve God more diligently. But let us now see what is the treatise of Saint John concerning the same.

First, he demonstrates that all that follows is not earthly, but heavenly. For he sees another token in heaven. He says “another,” because in chapter 12 we heard that mention was made of another certain sign. And he calls that a sign or token, which signifies another thing, and therefore is not to be considered of itself, but in as much as it brings into knowledge another certain thing, and that much greater than it shows at the first sight. He calls this sign, that is to wit, that same vision, great and marvelous. For the judgments of God are greatest and most wonderful. While they are executed, the ungodly marvel, which had thought such things should never have come to pass; the godly also marvel at the great power of God, his most just righteousness, and his ripeness and faithfulness in delivering and saving his people. Then he declares what sign was shown him in heaven, and by that celestial vision: he saw seven Angels, having in seven cups, plagues. That is, he perceived God prepared and furnished with divine power, with which he both might and would send plagues and deserved punishments, as well upon Antichrist himself, as upon his members, and all the ungodly men on Earth, for their wickedness committed against God. And as we have many times warned you in this book, the seventh number is the measure of fullness. Wherefore God has ministers enough and enough, by whose service he may plague and destroy the ungodly. And therefore seven plagues are all manner of plagues. Temporal plagues are abundantly recited in Leviticus 26 and Deuteronomy 28. The Lord is rich, and in everlasting plagues of most diverse kinds also. For Scripture in certain places rehearses a gnawing worm, a fire unquenchable, weeping and gnashing of teeth, outward darkness, and many others of like sort. But these seven plagues he calls the last, and immediately shows the reason, for in them is the wrath of God fulfilled. For on those last and most corrupt ages the Lord will pour out his plague, and that the plagues of his just wrath, and shall pour them out most fully to the end, and shall execute his full wrath against the ungodly, forevermore.

Yet now he suspends a while that narration concerning the Anglo-Saxon rulers of the plagues and places or sends before the great joys of the blessed Martyrs, triumphs, songs of praise, rejoicing, and thanksgiving. And this joy is interlaced here in the treatise of punishments, for the consolation of the faithful, that they should know themselves delivered from punishments. And if it happens, while the wicked are punished, that some displeasure touches them also (as it cannot be avoided, but that when the wicked are plagued, some detriment must also arise for the faithful), that they may understand yet that the dangers and detriments must be recompensed with the excellent abundance of joys. For hereby is signified how the godly rejoice while the Lord executes his justice. To be also the changeable course of things, that those who have once wept in the world, should now be glad and joyful, according to the saying of our Savior in John 16. Moreover, it behooved by the testimony of all Saints to be declared to the Saints who dwell on Earth, that the judgments of God are righteous and true; which thing, understood, questions and sundry murmurings against God cease.

\index{Arethas of Caesarea}\index{Daniel (Book of)}First he sees those who overcame Antichrist, and have had nothing to do with him; as we say in Dutch, for this I suppose is signified by that plentiful rehearsal of certain members (the declaration of which is set forth before) in heaven, not in some torturous place, or nowhere, as some men gather. He saw, I say, in heaven the blessed souls standing on a glassy sea, mixed with fire. And in another place I have told you that the sea figures the world, because of its rage and instability. Certainly Daniel takes it so in chapter 7. And it is called glassy because of its frailty and brittleness. For worldly things shine, but they are soon broken. Whereupon it is said that worldly things are as brittle as glass: which, while they shine, break. And not without cause is fire mixed with worldly things. For the saints, while they are conversant in earth, always feel in a manner the fire of affliction, as spoke of by Saint Peter, 1 Peter 4. And they stand on a glassy sea mingled with fire. For conquerors tread upon the world, and upon all the torments and mockeries of the world, triumphing over all worldly things. The prophet in Psalm 66 brings in the saints singing a joyful song unto God, and among other things saying, “You have brought us into snares, you have laid tribulations upon our backs, you have set men in our necks. We have passed through fire and water, and you have brought us out into a place of relief.” Therefore, alterations follow in another world. Wherefore Arethas, expounding this place, says that the glassy sea seems to indicate nothing else than that by the sea is meant the multitude, by the glass the brightness, by fire the purity of those who are worthy of that blessed life. And certainly the same words in diverse respects may signify diverse things, and make the sense agreeable.

Hitherto we have heard that the saints are in heaven, where they triumph over the world, which is vanquished. But now we shall hear more clearly what they do in heaven, and how they sing to the Lord a song of thanks and praises, which fully agrees with Psalm 66. And he attributes harps to the blessed martyrs, as he did to the Elders. These he calls of God, as you would say divine and celestial, fitting to set forth the praises of God. For a celestial jubilee is signified, of which it is spoken in the fifth chapter. He adds moreover, to express the music: and they sing. And he declares also the manner of their singing, the song of Moses, the servant of God, and the song of the Lamb. Therefore, this song of the saints is a rejoicing ditty, triumphant and of thanksgiving. For just as in times past Mary with the company of Israelite virgins, at the appointment of Moses, sang a song when the Israelites were delivered out of the bondage of Egypt, and Pharaoh was drowned in the Red Sea with his whole army—of which you may read more in the fourteenth and fifteenth chapters of Exodus—so the blessed souls in heaven praise God, who has delivered them from Satan, Antichrist, and the world. And the song of the Lamb is the Christian thanksgiving, by which the virtue of Christ and his redemption is praised by the saints. For just as the old fathers after eating the Paschal lamb made a jubilee and gave God thanks, so the blessed saints now enfranchised with the full liberty of the children of God, give thanks unto Christ their deliverer.

Finally, they recite the order and form of their song. God is highly commended here, called the Lord, God, almighty, King of Saints—the one for whom the saints fight, by whom they are governed, and who defends, maintains, and keeps the saints. He is called holy, in whom there is no blemish, no iniquity. And above all, they praise his works, which they call great and marvelous. These are manifest in heaven and on earth. They declare the power, wisdom, and justice of God. Therefore, they immediately conclude that God’s ways—that is, the considerations he follows in governing and doing things—are true and just. For he deceives not, he does no man wrong. Therefore, God is just in punishing the Antichrist and delivering his own. For although he may seem to neglect them, he keeps faith with the godly, just as a king never neglects his subjects.

Now they claim that it is fitting for all people on Earth to do likewise: it is right that all men fear you and glorify you in all things, neither to accuse nor murmur at your judgments. There is added another reason: he alone is holy, without sin, and without spot. None of all creation possesses this. Although many Gentiles now despise God, yet they will come and worship; they will know their own filthiness and the holiness and righteousness of God. For the justice and judgment of God, which are not yet revealed and therefore are despised, shall be revealed, that all the godly of all nations may attribute glory to the righteous God. These things also prepare the reader and hearer for the treatise now following concerning the judgments of God and the punishments of the ungodly. The Lord open the eyes of our minds, that we may see these things with fruitful understanding.
